\begin{lemma}[A finite shift orbit is decided by one finite prefix]
Let $h$ be locally constant, let $x$ be an increasing enumeration, and let
$m\in\mathbb N$.  There is a finite prefix of $x$ which determines all of
the finitely many values
\[
h(x),h(\mathsf{ShiftMap}(x)),\ldots,
h(\mathsf{PowerShiftIter}(m,x)).
\]
Precisely, there is $N$ such that whenever
$\mathsf{PrefixAgree}(N,x,y)$, then for every $j\leq m$,
\[
h(\mathsf{PowerShiftIter}(j,y))
=h(\mathsf{PowerShiftIter}(j,x)).
\]
\end{lemma}
\begin{proof}
Fix $j\leq m$.  By the definition of local constancy
\noderef{LocallyConstant}, local constancy of $h$ at
$\mathsf{PowerShiftIter}(j,x)$ supplies a number $n_j$ such that any
enumeration agreeing with this shifted enumeration through $n_j$ entries has
the same $h$-value.  There are only finitely many indices $j\leq m$.
Choose $N$ at least $n_j+j$ for each of them.  If
$\mathsf{PrefixAgree}(N,x,y)$, then
\noderef{PowerShiftIterPrefixAgree} gives
\mathsf{PrefixAgree}(n_j,\mathsf{PowerShiftIter}(j,x),
\mathsf{PowerShiftIter}(j,y)); local constancy gives the required equality.
The same $N$ works for all $j\leq m$, which proves the lemma.
\end{proof}
